% Shared notation; copied from docs/paper1/macros.tex at the split (S1.5).
\newcommand{\DKY}{D_{\mathrm{KY}}}
\newcommand{\dlat}{d_{\mathrm{lat}}}
\newcommand{\Nens}{N_{\mathrm{ens}}}
\newcommand{\R}{\mathbb{R}}
\newcommand{\Z}{\mathbb{Z}}
\newcommand{\E}{\mathcal{E}}          % encoder
\newcommand{\Dec}{\mathcal{D}}        % decoder
\newcommand{\Fp}{\mathcal{F}}         % latent propagator
\newcommand{\Hobs}{\mathcal{H}}       % physical observation operator
\newcommand{\Lat}{\Lambda}            % site lattice
\newcommand{\Shift}{\mathsf{S}}       % lattice shift
\newcommand{\Tr}{\mathsf{T}}          % physical translation
\newcommand{\Exp}{\mathbb{E}}
\newcommand{\Var}{\operatorname{Var}}
\newcommand{\Cov}{\operatorname{Cov}}
\newcommand{\tr}{\operatorname{tr}}
\newcommand{\dist}{\operatorname{dist}}
\newcommand{\Renc}{R_{\mathrm{enc}}}
\newcommand{\Rdec}{R_{\mathrm{dec}}}
\newcommand{\Rprop}{R_{\mathrm{prop}}}
\newcommand{\kmax}{k_{\max}}
\newcommand{\kmid}{k_{\mathrm{mid}}}
\newcommand{\wvar}{w_{\mathrm{var}}}
\newcommand{\wld}{w_{\mathrm{ld}}}
\newcommand{\ie}{i.e.\@\xspace}
\newcommand{\eg}{e.g.\@\xspace}

\newtheorem{proposition}{Proposition}
\newtheorem{lemma}{Lemma}
\newtheorem{corollary}{Corollary}
\theoremstyle{remark}
\newtheorem{remark}{Remark}

% Run-in paragraph heading (revtex letters \paragraph headings; we do not want that).
\newcommand{\para}[1]{\par\medskip\noindent\textbf{#1}\enspace}

% Wide table: full line width with the columns spread out (like revtex's ruledtabular) when it fits,
% shrunk to the line width only when its natural width is larger -- never enlarged, so a narrow table
% keeps the float font size (fixed 2026-10-11: \resizebox alone scaled small tables up).
% #1: column spec, #2: body (header + rows). The body is typeset twice (measured, then set).
\makeatletter
\newsavebox\wt@box
\newcommand{\widetabular}[2]{%
  \sbox\wt@box{\begin{tabular}{#1}\hline\hline\noalign{\vskip 2pt}#2\noalign{\vskip 2pt}\hline\hline\end{tabular}}%
  \ifdim\wd\wt@box>\linewidth
    \resizebox{\linewidth}{!}{\usebox\wt@box}%
  \else
    \begin{tabular*}{\linewidth}{@{\extracolsep{\fill}}#1}\hline\hline\noalign{\vskip 2pt}#2\noalign{\vskip 2pt}\hline\hline\end{tabular*}%
  \fi}
\makeatother
% Table bodies are read with TeX's primitive \input (expandable), so a rule may follow them.
\makeatletter\newcommand{\tabinput}[1]{\@@input #1.tex }\makeatother

% REVTeX's ruledtabular (used by every table) for non-REVTeX classes: centred table with
% double rules above and below; the table body supplies its own \hline after the header.
\makeatletter
\@ifundefined{ruledtabular}{%
\newenvironment{ruledtabular}{\par\centering
  \let\rt@tabular\tabular \let\rt@endtabular\endtabular
  \renewenvironment{tabular}[1]{\rt@tabular{##1}\hline\hline\noalign{\vskip2pt}}%
    {\noalign{\vskip2pt}\hline\hline\rt@endtabular}}{\par}}{}
\makeatother

% Placeholder for numbers that pending runs will fill (Block K; see docs/paper1/partII_narrative.md,
% table P1-P12). Renders in red so none survives into a submitted draft unnoticed.
\providecommand{\textcolor}[2]{#2}
\newcommand{\placeholder}[1]{\textcolor{red}{\textbf{[#1]}}}
